\chapter{Pointers}

\textbf{Readings:} CP:AMA 11, 17.7.

\tikzset{
  cell/.style={draw,minimum width=1.4cm,minimum height=0.7cm},
  ptr/.style={-{Stealth[length=2mm]},thick}
}

\section{The Address Operator}

The \emph{address operator} \code{\&} produces the (starting) address of where the value of an
identifier is stored in memory. The \code{printf} placeholder for an address is \code{\%p}
(printed in hexadecimal).
\begin{lstlisting}[style=c]
int g = 42;
int main(void) {
  printf("the value of g is %d\n", g);
  printf("the address of g is %p\n", &g);   // e.g. 0x601040
}
\end{lstlisting}

\section{Pointers}

\begin{definition}
A \emph{pointer} is a value that is an address. A pointer variable is defined by placing a star
\code{*} before the identifier.
\end{definition}

\begin{lstlisting}[style=c]
int i = 5;
int *p = &i;    // p "points at" i
\end{lstlisting}

\begin{center}
\begin{tikzpicture}
  \node[cell] (i) {5};
  \node[left=1mm of i] {\code{i}};
  \node[below=0mm of i,font=\scriptsize] {0x5000};
  \node[cell,right=3cm of i] (p) {0x5000};
  \node[right=1mm of p] {\code{p}};
  \node[below=0mm of p,font=\scriptsize] {0x5008};
  \draw[ptr] (p.west) -- (i.east);
\end{tikzpicture}
\end{center}

The type of \code{p} is ``pointer to int'', written \code{int *}. Pointers have types:
an \code{int *} should only point at \code{int}s, a \code{char *} at \code{char}s, and so on.

\begin{aside}
The placement of the space does not matter: \code{int *p}, \code{int* p} and \code{int * p} are
all the same. CS 136 style puts the star next to the identifier because
\code{int *p, q;} defines a pointer \code{p} and an \emph{int} \code{q}. Avoid defining several
pointers on one line.
\end{aside}

\subsection{Size of a pointer}

On a $k$-bit system, addresses (and hence all pointers, regardless of the type they point to)
are $k$ bits. Seashell is a 64-bit system, so \code{sizeof(int *)} is \textbf{8 bytes}.

\section{The Indirection Operator}

The \emph{indirection} operator \code{*} (also called the \emph{dereference} operator) is the
inverse of \code{\&}: \code{*p} is ``the value at the address \code{p}''.
\begin{lstlisting}[style=c]
int i = 5;
int *p = &i;
printf("%d\n", *p);   // 5
*p = 7;               // changes i!
printf("%d\n", i);    // 7
\end{lstlisting}
\code{*\&i} is the same as \code{i}. Note that the star means two different things: in a
definition (\code{int *p}) it is part of the type, and in an expression (\code{*p}) it is the
indirection operator.

Pointers can point to pointers:
\begin{lstlisting}[style=c]
int i = 5;
int *p = &i;
int **pp = &p;     // pointer to pointer to int
**pp = 6;          // i is now 6
\end{lstlisting}

\begin{center}
\begin{tikzpicture}
  \node[cell] (pp) at (0,0) {0x5008}; \node[above=0mm of pp] {\code{pp}};
  \node[below=0mm of pp,font=\scriptsize] {0x5010};
  \node[cell] (p) at (3.4,0) {0x5000}; \node[above=0mm of p] {\code{p}};
  \node[below=0mm of p,font=\scriptsize] {0x5008};
  \node[cell] (i) at (6.8,0) {6}; \node[above=0mm of i] {\code{i}};
  \node[below=0mm of i,font=\scriptsize] {0x5000};
  \draw[ptr] (pp.east) -- (p.west);
  \draw[ptr] (p.east) -- (i.west);
  \node[font=\scriptsize,anchor=west,align=left] at (7.7,0) {\code{*pp} is \code{p}\\\code{**pp} is \code{i}};
\end{tikzpicture}
\end{center}

\section{The \code{NULL} Pointer}

\code{NULL} (defined in \code{<stdlib.h>} and \code{<stdio.h>}, among others) is a special
pointer value that points to nothing; it is address 0. It is used to represent ``no value''
(like \code{empty} or \code{\#f} in Racket).

\begin{important}[Dereferencing \code{NULL}]
Dereferencing a \code{NULL} pointer (or any invalid pointer) is undefined behaviour; in practice
the program crashes with a segmentation fault. Since \code{NULL} is false and every other
pointer is true, \code{assert(p);} asserts that \code{p} is not \code{NULL}. Functions that take
pointers should usually start with such asserts (and the documentation should say
``requires: p is a valid pointer'').
\end{important}

\section{Pointer Assignment vs.\ Value Assignment}

This is the single most important idea in the unit.
\begin{lstlisting}[style=c]
int i = 5;
int j = 6;
int *p = &i;
int *q = &j;
p = q;        // pointer assignment
\end{lstlisting}
\code{p = q} changes \emph{what \code{p} points at}: now both \code{p} and \code{q} point at
\code{j}, so \code{*p} is 6. \code{i} is still 5.
\begin{center}
\begin{tikzpicture}
  \node[cell] (i) {5}; \node[left=1mm of i] {\code{i}};
  \node[cell,below=8mm of i] (j) {6}; \node[left=1mm of j] {\code{j}};
  \node[cell,right=3cm of i] (p) {}; \node[right=1mm of p] {\code{p}};
  \node[cell,right=3cm of j] (q) {}; \node[right=1mm of q] {\code{q}};
  \draw[ptr] (p.center) -- (j.north east);
  \draw[ptr] (q.center) -- (j.east);
  \draw[ptr,dashed,gray] (p.west) -- node[above,font=\scriptsize]{before} (i.east);
\end{tikzpicture}
\end{center}

\begin{lstlisting}[style=c]
int i = 5;
int j = 6;
int *p = &i;
int *q = &j;
*p = *q;      // value assignment
\end{lstlisting}
\code{*p = *q} copies the \emph{value} that \code{q} points at into the location \code{p} points
at. \code{p} still points at \code{i}, but now \code{i} is 6.

\begin{definition}
When two (or more) pointers point to the same location, or a pointer points to a named
variable, we say they are \emph{aliases}: the same memory can be changed through either name.
After \code{p = q} above, \code{*p} and \code{*q} are aliases for \code{j}.
\end{definition}

Aliasing makes code harder to reason about: changing \code{*p} silently changes \code{*q}.

\section{Mutation Through Parameters}

Since C passes arguments by value, a function cannot change the caller's variables directly.
But if we pass the \emph{address} of a variable, the function can change the variable through
the pointer. This emulates \emph{pass by reference}.
\begin{lstlisting}[style=c]
// inc(p) increments the integer p points at
// requires: p is a valid pointer
// effects: modifies *p
void inc(int *p) {
  assert(p);
  *p += 1;
}

int main(void) {
  int x = 5;
  inc(&x);              // x is now 6
}
\end{lstlisting}

\subsection{The classic example: swap}
\begin{lstlisting}[style=c]
// swap(px, py) swaps the values of *px and *py
// effects: modifies *px and *py
void swap(int *px, int *py) {
  int temp = *px;
  *px = *py;
  *py = temp;
}

int a = 3, b = 4;
swap(&a, &b);    // a is 4, b is 3
\end{lstlisting}
A version \code{void swap(int x, int y)} would only swap the copies in its own stack frame and
have no effect on the caller.

\begin{center}
\begin{tikzpicture}[row/.style={minimum width=4.4cm,minimum height=0.55cm,text width=4.2cm,align=left,inner sep=2pt}]
  \node[row,font=\bfseries] (sw) {swap};
  \node[row,below=0pt of sw] (px) {px:};
  \node[row,below=0pt of px] (py) {py:};
  \node[row,below=0pt of py] (tmp) {temp: 3};
  \node[row,below=0pt of tmp] (ra) {return address: main};
  \node[row,below=0pt of ra,font=\bfseries] (mn) {main};
  \node[row,below=0pt of mn] (a) {a: 3};
  \node[row,below=0pt of a] (b) {b: 4};
  \node[row,below=0pt of b] (rb) {return address: OS};
  \draw (sw.north west) rectangle (ra.south east);
  \draw (mn.north west) rectangle (rb.south east);
  \draw[ptr] ($(px.west)+(0.9,0)$) -- ($(px.east)+(0.9,0)$) |- (a.east);
  \draw[ptr] ($(py.west)+(0.9,0)$) -- ($(py.east)+(0.5,0)$) |- (b.east);
\end{tikzpicture}
\end{center}

\subsection{Returning more than one value}

A function can only \code{return} one value, but it can ``produce'' more by storing results
through pointer parameters (\emph{output parameters}):
\begin{lstlisting}[style=c]
// divide(num, denom, quot, rem) stores num / denom in *quot
//   and num % denom in *rem
// requires: denom != 0, quot and rem are valid pointers
// effects: modifies *quot and *rem
void divide(int num, int denom, int *quot, int *rem) {
  assert(denom != 0);
  assert(quot && rem);
  *quot = num / denom;
  *rem = num % denom;
}
\end{lstlisting}
A common pattern is to return a \code{bool} success flag and store the actual result through a
pointer (this is exactly how \code{scanf} works, Unit~07).

\begin{important}[Never return the address of a local variable]
\begin{lstlisting}[style=c]
int *bad(void) {
  int x = 5;
  return &x;     // x's frame is popped on return: dangling pointer!
}
\end{lstlisting}
Once the function returns, its stack frame is gone; the pointer now points at memory that will
be reused by the next call. Using it is undefined behaviour.
\end{important}

\begin{center}
\begin{tikzpicture}[fr/.style={draw,minimum width=4.2cm,align=left,text width=4cm,inner sep=4pt,font=\small}]
  \node[font=\scriptsize\bfseries] at (0,2.3) {during \code{bad()}};
  \node[fr] (b) at (0,1.35) {\textbf{bad}\\ x: 5};
  \node[fr,below=0pt of b] (m) {\textbf{main}\\ p: ???};
  \node[font=\scriptsize\bfseries] at (6,2.3) {after \code{p = bad();}};
  \node[fr,draw=gray,dashed,text=gray] (b2) at (6,1.35) {(popped frame)\\ x: 5 \emph{(stale)}};
  \node[fr,below=0pt of b2] (m2) {\textbf{main}\\ p:};
  \coordinate (d) at ($(m2.south west)+(0.75,0.3)$);
  \fill (d) circle (1.5pt);
  \draw[ptr,red!70!black] (d) -- ($(m2.east |- d)+(0.4,0)$) |- (b2.east);
  \node[font=\scriptsize,red!70!black,anchor=west] at ($(m2.east)+(0.5,0.45)$) {dangling};
\end{tikzpicture}
\end{center}
The next function call reuses the popped frame's memory, overwriting whatever \code{p} points at.

\section{Passing Structures}

Passing a structure by value copies the entire structure into the new stack frame. For large
structures that wastes time and stack space (and can contribute to stack overflow). Instead we
pass a pointer, which is always 8 bytes in Seashell regardless of the structure size.

Fields are accessed through a pointer with \code{(*ptr).field}. The parentheses are needed
because \code{.} has higher precedence than \code{*}. This is so common that C provides the
\emph{arrow operator}, which combines indirection and selection: \code{ptr->field}.
\begin{lstlisting}[style=c]
struct posn { int x; int y; };

// posn_scale(p, k) multiplies both coordinates of *p by k
// effects: modifies *p
void posn_scale(struct posn *p, int k) {
  assert(p);
  p->x *= k;           // same as (*p).x *= k
  p->y *= k;
}
\end{lstlisting}

\subsection{\code{const} pointers}

If a function takes a pointer only for efficiency and does not modify the data, mark it
\code{const} so the compiler enforces (and documents) that:
\begin{lstlisting}[style=c]
int posn_dist_sq(const struct posn *a, const struct posn *b) {
  int dx = a->x - b->x;
  int dy = a->y - b->y;
  return dx * dx + dy * dy;
}
\end{lstlisting}
The position of \code{const} matters:
\begin{lstlisting}[style=c]
const int *p;          // p can change; *p cannot be changed through p
int * const q = &i;    // q cannot change; *q can
const int * const r = &i;  // neither can change
\end{lstlisting}
(Read declarations right to left: ``\code{q} is a const pointer to an int''.)

\section{Opaque Structures}

C allows a structure to be \emph{declared} without being \emph{defined}:
\begin{lstlisting}[style=c]
struct account;     // incomplete type: fields unknown
\end{lstlisting}
A client that only sees this declaration cannot create a \code{struct account} or access its
fields (it does not know its size or fields), but it \emph{can} use pointers to it
(\code{struct account *}), since all pointers have the same size. This is an \emph{opaque
structure}, and it is how C enforces information hiding for ADTs: the header declares the
structure, the implementation file defines it, and clients can only use the interface functions.
\begin{lstlisting}[style=c]
// account.h
struct account;
struct account *account_create(int id, int balance);   // Unit 10
int account_balance(const struct account *acc);
void account_deposit(struct account *acc, int amount);
\end{lstlisting}
\begin{lstlisting}[style=c]
// account.c
#include "account.h"
struct account {       // the complete definition is private to this file
  int id;
  int balance;
};
int account_balance(const struct account *acc) {
  return acc->balance;
}
\end{lstlisting}
Creating opaque structures requires dynamic memory, which we see in Unit~10.

\section{Function Pointers}

In Racket, functions are first-class values. In C, functions are not values, but \emph{pointers
to functions} are: the name of a function evaluates to the address of its code (in the code
section). This lets us pass a function to another function.
\begin{lstlisting}[style=c]
int my_sqr(int n) { return n * n; }
int my_neg(int n) { return -n; }

// a function pointer variable: f points to a function int -> int
int (*f)(int) = my_sqr;
printf("%d\n", f(5));     // 25
f = my_neg;
printf("%d\n", f(5));     // -5

// apply_twice(f, x) produces f(f(x))
int apply_twice(int (*f)(int), int x) {
  return f(f(x));
}
apply_twice(my_sqr, 3);   // 81
\end{lstlisting}
\begin{center}
\begin{tikzpicture}[fn/.style={draw,rounded corners,fill=gray!8,font=\scriptsize\ttfamily,minimum width=2.4cm,minimum height=7mm}]
  \node[font=\scriptsize\bfseries] at (5,1.15) {code section};
  \draw[gray] (3.4,-1.75) rectangle (6.6,0.85);
  \node[fn] (sq) at (5,0.25) {my\_sqr: ...};
  \node[fn] (ng) at (5,-1.05) {my\_neg: ...};
  \node[cell] (f) at (0,-0.4) {}; \node[left=1mm of f] {\code{f}};
  \fill (f.center) circle (1.5pt);
  \draw[ptr,dashed,gray] (f.center) -- node[above,font=\scriptsize]{before} (sq.west);
  \draw[ptr] (f.center) -- node[below,font=\scriptsize]{after \code{f = my\_neg}} (ng.west);
\end{tikzpicture}
\end{center}
The syntax \code{int (*f)(int)} reads: ``\code{f} is a pointer to a function taking an
\code{int} and returning an \code{int}''. The parentheses around \code{*f} are needed:
\code{int *f(int)} would declare a function returning an \code{int *}. We use function pointers
for ``map''-style array functions (Unit~08) and generic ADTs (Unit~13).

\section{Common Pitfalls}
\begin{itemize}
  \item Dereferencing an uninitialized or \code{NULL} pointer (segmentation fault or worse).
  \item Returning the address of a local variable (dangling pointer).
  \item Confusing \code{p = q} (pointer assignment) with \code{*p = *q} (value assignment).
  \item Writing \code{*p.x} instead of \code{p->x}.
  \item Forgetting \code{\&} when passing a variable that should be modified
    (e.g.\ \code{scanf("\%d", n)}).
\end{itemize}

\section{Summary}
\begin{itemize}
  \item \code{\&x} is the address of \code{x}; \code{*p} is the value at address \code{p}.
  \item Pointers are typed; all pointers are 8 bytes in Seashell; \code{NULL} points nowhere.
  \item Pass a pointer to let a function modify the caller's data (emulated pass by reference);
    use output parameters to produce several results.
  \item Pass structures by pointer; use \code{->}; use \code{const} for read-only access.
  \item Opaque structures (declared but not defined in the header) enforce ADT abstraction.
  \item Function pointers let functions be passed as arguments.
\end{itemize}
